\section{The Other Half of a Dispatch}
\label{sec:ch16-when-a-batch-goes-out}
\index{DataLoader!the executor's half|(}

Chapter~\ref{ch:data-without-the-n-plus-1} read this machinery already, and got
to a sentence it could not finish. It found that the documentation describes the
dispatch trigger as \enquote{no more resolver work is immediately ready}, and
answered:

\begin{quote}
\enquote{Nothing in the code that dispatches batches observes resolver
readiness. It observes the batch.}
\end{quote}

That is correct, and it is only half of an arrangement that has two halves.
Something in this system does observe resolver readiness. It is the scheduler
from the previous section, and what it does with the observation is not what the
documentation implies.

\index{WorkScheduler@\code{WorkScheduler}!waking the dispatcher}%
The scheduler runs \code{TryDispatchOrCompleteUnsafe} whenever the drive loop
comes round, in \code{WorkScheduler.Execute.cs}, and this is the part that
matters:

\begin{minted}[fontsize=\footnotesize,breaklines]{csharp}
if (!isWaitingForTaskCompletion)
{
    isWaitingForTaskCompletion = _work is { HasRunningTasks: true, IsEmpty: true };
}

var hasWork = !_work.IsEmpty || !_serial.IsEmpty;

if (isWaitingForTaskCompletion)
{
    _signal.Reset();

    if (Interlocked.CompareExchange(ref _hasBatches, 0, 1) == 1)
    {
        _batchDispatcher.BeginDispatch(_ct);
    }
}
\end{minted}

\index{WorkScheduler@\code{WorkScheduler}!what idle means here}%
The condition is \code{HasRunningTasks: true, IsEmpty: true}: there is nothing
left to hand out, and something is still in flight. That is what \enquote{every
resolver that can run is now blocked on something} looks like from inside, and
it is a fair reading of resolver readiness.

But \code{BeginDispatch} does not dispatch anything. It sets a signal and makes
sure the coordinator loop is running. The coordinator then applies the rule
chapter~\ref{ch:data-without-the-n-plus-1} printed, which asks the batch about
itself and never asks the scheduler anything. So both sentences are true at
once: readiness is observed, and it is not what decides. It decides when to
start asking.

\index{DataLoader!the seam between the two halves}%
Which finally locates the gap that chapter's two runs in four hundred came out
of. It is not inside the dispatcher and not inside the scheduler. It is between
them: the scheduler knows when work has stalled, the batch knows when it last
grew, and no line anywhere relates either fact to how many keys the batch is
still waiting for. Nothing in the system holds that number. The settle window is
what stands in for it, and a straggler that misses the window is what a stand-in
costs.

\index{DataLoader!what is not being claimed}%
That is as far as I will take it, because it would be easy to write a satisfying
sentence here that the evidence does not carry. The source shows the settle rule
is the only thing that can release a batch, so it is the only thing that can
have released one early. It cannot tell me why a particular key was late on
those two runs. Thread-pool scheduling is the obvious candidate and I have not
measured it, so a candidate is what it stays. The experiment that would settle
it is to set \code{BatchSettleTimeUs} to zero, which the option documents as
making a batch eligible as soon as two evaluations see it unchanged, and run the
four hundred requests again. I did not, and the note says so.

\subsection*{Two things called batching}

\index{DataLoader!not the same as a batch resolver}%
One distinction before leaving this, because two unrelated mechanisms here share
a word and only one of them is the one you know.

\code{Batch} and \code{BatchDispatcher} are GreenDonut's DataLoader key
batching, which is everything above and everything
chapter~\ref{ch:data-without-the-n-plus-1} measured.
\code{BatchResolverTask} is HotChocolate's own thing entirely: a grouping of
resolver invocations that share one field-selection path, created when a
selection's strategy is batch, and released when its ancestor paths have
drained rather than on any timer. It is the fourth of the four strategies
section~\ref{sec:ch16-pure-or-a-task} listed and the one Mosaic never uses.

They meet in the scheduler and share nothing else, including their dispatch
conditions. A profiler that shows you both will name them almost identically.

\subsection*{And how it all ends}

\index{WorkScheduler@\code{WorkScheduler}!knowing the operation is done}%
The same method decides the operation is finished, in its other branch: when
neither queue has work and no batch entries are pending, \code{\_isCompleted}
goes true and the drive loop falls out. A resolver blocked on a DataLoader
promise does not count as absent while it waits, because the queue's running
counter is incremented when the task is taken and decremented when it completes,
so the operation cannot be declared done underneath it.

\code{AsyncManualResetEvent} is the primitive holding all of this together, and
it is hand-written rather than borrowed: an awaitable of seventy-four lines, in
a file of its own, which resumes its continuation by queueing it to the thread
pool. There is a \code{SemaphoreSlim} shaped hole in this design, and somebody
deliberately did not put one in it.

\medskip
So much for the machinery that runs a request. The last section is about three
middleware whose whole purpose is that a request does not get that far, and one
of them turns out to have been switched on by a chapter that was not thinking
about caching at all.

\index{DataLoader!the executor's half|)}
